\begin{tabular}{lrrrr}
\toprule
Scenario & Links & Median change (\%) & Minimum (\%) & Maximum (\%)\\
\midrule
2024 held constant (reference) & 45 & 0.00 & 0.00 & 0.00\\
Half of the 2012--2024 trend extended & 45 & $-3.36$ & $-8.53$ & $+6.06$\\
Full 2012--2024 trend extended & 45 & $-6.94$ & $-17.12$ & $+8.65$\\
\bottomrule
\end{tabular}
\\[0.5em]
\noindent\TODO{Final note, in your own words. Change is relative to the 2024 surface.
Only built-up area and vegetation (tree, non-tree, bare) were pushed forward along their
2012--2024 straight-line trend; nighttime lights, roads, livestock, terrain, fences, and
the cores were held at 2024. These are ``what if the recent trend kept going'' tests, not
forecasts.\\
WHY THE ROWS WERE RENAMED: they used to say ``Low growth,'' ``Continuation,'' and ``High
development.'' Those names were misleading. The ``low growth'' row is zero by
construction because nothing changed. The ``high development'' row extends the whole
trend, and that includes vegetation. Vegetation made links easier from 2012 to 2024, so
extending it makes links easier still ($-6.94\%$). A reader who sees ``high development''
expects connections to get harder, so the old label would look like a mistake.\\
WHAT TO SAY ABOUT IT: one sentence noting that the projected change is mostly driven by
the vegetation trend, not by new development. If you want a true ``more development''
test, it would need built-up and lights pushed forward with vegetation held at 2024;
that run does not exist yet.}
